\documentclass[10pt,a4paper]{amsart}
\usepackage[margin=19mm]{geometry}
\usepackage{amsmath,amssymb,mathtools}
\usepackage[scr=rsfs,cal=euler]{mathalfa}
\usepackage{xfrac}
\usepackage[unicode,hidelinks,bookmarks=false]{hyperref}
\newcommand{\ie}{i.e.\ }
\newcommand{\cf}{cf.\ }
\newcommand{\resp}{respectively\ }
\newcommand{\Subsection}{\subsection}
\providecommand{\nobreakdash}{\nobreak\hbox{-}}
\setlength{\emergencystretch}{2em}
\multlinegap0pt
\usepackage{amssymb}
\usepackage{mathtools}
\usepackage[shortlabels]{enumitem}
\newtheorem{lemma}{Lemma}
\newcommand{\Sym}{\operatorname{Sym}}
\title{Free-Energy Asymptotics of the Two-Dimensional Quantum Heisenberg Ferromagnet}
\author{Andreas Klippel}
\date{Source: arXiv:2608.25506v1 (CC BY 4.0)}
\begin{document}
\maketitle

\noindent\emph{Source and licence.} Free-Energy Asymptotics of the Two-Dimensional Quantum Heisenberg Ferromagnet. Version arXiv:2608.25506v1 (CC BY 4.0), distributed by arXiv under the Creative Commons Attribution 4.0 International licence, http://creativecommons.org/licenses/by/4.0/, as recorded in the arXiv licence metadata for this exact version and shown on the abstract page https://arxiv.org/abs/2608.25506v1. Full licence notice: \texttt{licenses/arxiv-2608.25506-CC-BY-4.0.txt}.\par\smallskip

\noindent\textit{Editorial scope.} Counted unit: the article's lemma lem:single-diagonal-trace (source0.tex lines 459--470). The article's complete original proof of it is delivered with it (source0.tex lines 472--661, one \texttt{proof} environment of 190 lines). No mathematics is written, repaired or re-derived: the statement and the proof are the article's own lines, in the article's own notation, and no retained line is substituted or re-flowed.  No further source-local context is needed: every cross-reference of every retained line resolves inside the two delivered spans.  Nothing was omitted from any retained span, so the package carries no omission record.  No retained line contains a citation, so the wrapper carries no bibliography and no bibliography entry.  No retained line contains a figure environment, an image-inclusion command or drawing code, so no image is delivered and the package declares no asset dependency.  Editorial formatting only, in addition: an AMS \texttt{amsart} preamble that restates the article's own declarations and macros used by the retained lines, namely the environments \texttt{lemma} on separate counters, together with the standard packages those lines need.  The retained environments print in this document as Lemma 1 in order of appearance, whereas the article prints the same items with the numbering of its own full document.  The complete native source of the article as distributed in the arXiv e-print for this exact version is bundled unchanged as \texttt{sources/208583/source0.tex}.
\begin{lemma}
\label{lem:single-diagonal-trace}
There is a universal \(C<\infty\) such that, for every periodic square with
integer \(L\geq3\), every integer \(k\geq2\), every
\(1\leq i<j\leq k\), and every symmetric
\(F\in\ell^2_{\Sym}(B_L^k)\),
\begin{equation}
 \sum_{e\in\mathcal B_{ij}}|\nabla_eF|^2
 \leq C\|L_{0,L,k}F\|^2.
 \label{eq:single-diagonal-trace}
\end{equation}
\end{lemma}

\begin{proof}
Set \(\mathbb Z_L:=\mathbb Z/L\mathbb Z\).  We use the unitary Fourier
transform
\[
 \widehat F(p_1,\ldots,p_k)
 :=V^{-k/2}\sum_{x_1,\ldots,x_k\in B_L}
 e^{-\mathrm i\sum_jp_j\cdot x_j}F(x_1,\ldots,x_k).
\]
Let
\[
 \mathbb T_L^2:=\frac{2\pi}{L}\mathbb Z_L^2,
 \qquad
 \varepsilon(p)=2\sum_{\nu=1}^2(1-\cos p_\nu).
\]
Fix the total momentum \(p\) of particles \(i,j\), their relative momentum
\(q\), and the momenta \(p_r\), \(r\ne i,j\), of the remaining particles,
which we call spectators.  Choose for the
class \(p\) its centred lift \(\widetilde p\in[-\pi,\pi)^2\); all
differences below are understood modulo \(2\pi\) with this lift.  On the
subspace with these momenta fixed, the product eigenvalue is
\begin{equation}
 \lambda(q):=\varepsilon(q)+\varepsilon(\widetilde p-q)+a,
 \qquad a:=\sum_{r\ne i,j}\varepsilon(p_r).
 \label{eq:product-fibre-energy}
\end{equation}
We abbreviate the resulting function of the remaining relative momentum
by \(\widehat F(q)\); the fixed values \(p,(p_r)_{r\ne i,j}\) are
suppressed.
The elementary identity
\[
 \varepsilon(q)+\varepsilon(\widetilde p-q)
 =4\sum_{\nu=1}^2
 \left[1-\cos(\widetilde p_\nu/2)
 \cos(q_\nu-\widetilde p_\nu/2)\right]
\]
implies, with constants independent of \(L,p,q\), that
\begin{equation}
 c\bigl\{\varepsilon(\widetilde p)
       +\varepsilon(q-\widetilde p/2)\bigr\}
 \leq \varepsilon(q)+\varepsilon(\widetilde p-q)
 \leq C\bigl\{\varepsilon(\widetilde p)
       +\varepsilon(q-\widetilde p/2)\bigr\}.
 \label{eq:pair-dispersion-comparison}
\end{equation}
Indeed, for one coordinate put
\(u=\cos(\widetilde p_\nu/2)\in[0,1]\) and
\(r=q_\nu-\widetilde p_\nu/2\).  Then
\[
 4(1-u\cos r)=4(1-u)+4u(1-\cos r),
 \qquad
 2(1-\cos\widetilde p_\nu)=4(1-u)(1+u).
\]
If \(u\geq1/2\), the second term in the first identity controls
\(1-\cos r\); if \(u<1/2\), the first term is bounded below by \(2\)
and hence controls the bounded quantity \(1-\cos r\).  The upper bound
follows from the same identities.

We shall use the following shifted lattice-sum estimate.  For every fixed
\(c_0>0\), uniformly in \(\theta\in\mathbb R^2\), \(L\geq3\), and
\(h\geq0\) satisfying either \(h=0\) or \(h\geq c_0L^{-2}\),
\begin{equation}
 \frac1V\sum_{q\in\mathbb T_L^2:\,h+\varepsilon(q-\theta)>0}
 \frac{h}{\bigl(h+\varepsilon(q-\theta)\bigr)^2}
 \leq C_{c_0}.
 \label{eq:shifted-lattice-resolvent}
\end{equation}
Here and below \(\varepsilon\) is evaluated modulo \(2\pi\).  To prove
\eqref{eq:shifted-lattice-resolvent}, let
\(d_{\mathbb T}(q,\theta)\) be Euclidean distance on the two-dimensional
torus.  On a centred fundamental domain,
\[
 c d_{\mathbb T}(q,\theta)^2
 \leq\varepsilon(q-\theta)
 \leq C d_{\mathbb T}(q,\theta)^2.
\]
Since \(\mathbb T_L^2\) has mesh \(2\pi/L\), a disk-packing argument gives,
uniformly in the shift,
\begin{equation}
 \#\{q\in\mathbb T_L^2:\varepsilon(q-\theta)\leq s\}
 \leq C(1+Vs),\qquad 0<s\leq8.
 \label{eq:shifted-grid-count}
\end{equation}
For completeness, place around every lattice point a disjoint square of
side \(\pi/L\).  All squares corresponding to momenta counted on the
left-hand side of \eqref{eq:shifted-grid-count} lie in a disk of radius
\(C\sqrt{s}+C/L\); comparison of areas gives
\eqref{eq:shifted-grid-count}.

If \(h=0\), every summand in \eqref{eq:shifted-lattice-resolvent} is zero.
If \(h\geq1\), every summand is at most \(h^{-1}\leq1\).  It remains to
consider \(c_0L^{-2}\leq h<1\).  Split the grid into
\[
 A_0=\{\varepsilon(q-\theta)\leq h\},\qquad
 A_j=\{2^{j-1}h<\varepsilon(q-\theta)\leq2^jh\},\quad j\geq1,
\]
stopping once \(2^jh\geq8\).  By \eqref{eq:shifted-grid-count}, the
contribution of \(A_0\) is at most
\[
 \frac{C(1+Vh)}{Vh}\leq C_{c_0},
\]
and the contribution of \(A_j\) is at most
\[
 \frac{C(1+V2^jh)}{V}\,
 \frac{h}{(2^{j-1}h)^2}
 \leq C_{c_0}4^{-j}+C2^{-j}.
\]
Summing in \(j\) proves \eqref{eq:shifted-lattice-resolvent}.  This proof
does not require \(\theta\) to be a lattice momentum; in particular it
covers the half-grid shift \(\theta=\widetilde p/2\).

Put
\[
 \mu:=a+\varepsilon(\widetilde p),\qquad
 \rho(q):=\varepsilon(q-\widetilde p/2).
\]
By \eqref{eq:pair-dispersion-comparison},
\(\lambda(q)\geq c(\mu+\rho(q))\).  Moreover, since \(p\) and all
spectator momenta lie in \(\mathbb T_L^2\), either \(\mu=0\) or
\(\mu\geq c_0L^{-2}\).  The smallest positive value of
\(\varepsilon\) on \(\mathbb T_L^2\) is
\(4\sin^2(\pi/L)\geq c_0L^{-2}\).

Restriction to \(x_i=x_j\) is the normalized sum
\(V^{-1/2}\sum_q\).  For any multiplier \(m(q)\), weighted
Cauchy--Schwarz gives
\begin{equation}
 \left|\frac1{\sqrt V}\sum_qm(q)\widehat F(q)\right|^2
 \leq
 \left(\frac1V\sum_{\lambda(q)>0}
 \frac{|m(q)|^2}{\lambda(q)^2}\right)
 \sum_q\lambda(q)^2|\widehat F(q)|^2.
 \label{eq:weighted-restriction-cs}
\end{equation}
Terms with \(\lambda(q)=0\) are omitted; the corresponding gradient
multiplier vanishes in every application below.

For tangential edges, Parseval in the spectator coordinate gives the sum
of squared multipliers
\[
 \sum_{r\ne i,j}\sum_{|e|=1}
 |e^{\mathrm i p_r\cdot e}-1|^2=2a,
\]
where using oriented rather than unoriented unit vectors changes only the
displayed factor.  Weighted Cauchy--Schwarz therefore reduces
the tangential contribution to
\begin{align}
 \frac1V\sum_{q:\lambda(q)>0}\frac{a}{\lambda(q)^2}
 &\leq \frac C V\sum_q\frac{a}{(\mu+\rho(q))^2} \notag\\
 &\leq \frac C V\sum_q\frac{\mu}{(\mu+\rho(q))^2}
 \leq C.
 \label{eq:tangential-restriction-sum}
\end{align}
If \(\mu=0\), then \(a=0\), so the left-hand side of
\eqref{eq:tangential-restriction-sum} vanishes and no \(0/0\) term is
present.

For a normal edge in a coordinate direction, permutation symmetry gives
\(\widehat F(q,\widetilde p-q)=
  \widehat F(\widetilde p-q,q)\).  Thus, after pairing the terms
\(q\) and \(\widetilde p-q\), the raw multiplier may be replaced exactly
by
\begin{equation}
 m_N(\widetilde p,q)
 :=\frac12\bigl(e^{\mathrm iq\cdot e}
 +e^{\mathrm i(\widetilde p-q)\cdot e}-2\bigr).
 \label{eq:normal-multiplier}
\end{equation}
With \(r=q-\widetilde p/2\), equation \eqref{eq:normal-multiplier} becomes
\(m_N=e^{\mathrm i\widetilde p_\nu/2}\cos r_\nu-1\).  Hence
\[
 |m_N|^2
 \leq C\bigl(\varepsilon(\widetilde p)+\rho(q)^2\bigr).
\]
Consequently,
\begin{align}
 &\frac1V\sum_{q:\lambda(q)>0}
 \frac{\varepsilon(\widetilde p)+\rho(q)^2}{\lambda(q)^2}
 \notag\\
 &\quad\leq
 \frac C V\sum_q\frac{\mu}{(\mu+\rho(q))^2}
 +\frac C V\sum_q\frac{\rho(q)^2}{(\mu+\rho(q))^2}
 \leq C.
 \label{eq:normal-restriction-sum}
\end{align}
The first sum is bounded by \eqref{eq:shifted-lattice-resolvent}; every
summand in the second is at most one.  When \(\mu=0\), the unique possible
zero of \(\lambda\) is omitted and its normal multiplier vanishes.
Summing the finitely many normal and tangential directions and using
Parseval proves \eqref{eq:single-diagonal-trace}.
\end{proof}

\end{document}
